日本血吸虫病所致肝纤维化是流行地区疾病管理中需要长期关注的器官损伤，可伴随门静脉高压及严重肝脾并发症 \citep{colley2014schistosomiasis,mcmanus2018schistosomiasis}。在包括中国在内的流行地区，感染控制与既有肝脏病变的评估共同构成疾病管理的重要环节 \citep{mcmanus2010china,lo2022who}。纵向研究表明，即使接受抗寄生虫治疗，严重的肝实质纤维化仍可能持续多年 \citep{carlton2010morbidity}。因此，能够重复开展、反映病变严重程度的影像评估，是连接人群筛查、病情评估与长期随访的重要基础。

B-mode 超声能够非侵入性地呈现肝实质与门脉相关改变，为上述评估提供了可重复使用的影像工具。然而，日本血吸虫病相关纤维化的超声表现具有复杂的空间形态：间隔性纤维化可呈现网状或鱼鳞样回声，门脉相关改变则具有不同的解剖分布；Basel 超声方案也强调对这些形态进行区分 \citep{richter2025basel}。这些特征使严重程度评估需要同时考虑局部回声、结构分布与解剖背景。对于自动分级系统，这提出了两个相互关联的要求：从复杂影像中提取足以支持细粒度评分的信息，以及将相关空间线索保留为能够与原图对照的输出。

既有计算研究已从超声纹理分析发展到放射组学和深度卷积模型，为减少人工判读差异提供了技术基础 \citep{yeh2003fibrosis,guo2024radiomics,lee2020ultrasound}。其中，SFibAI 针对日本血吸虫病相关肝纤维化，将超声评分表示为从 0.0 到 3.5、间隔为 0.1 的 36 个有序分值，并保留其与临床四级分级体系的对应关系 \citep{xu2026sfibai}。这一表示在粗粒度等级之外提供了更细致的严重程度描述，也为直接从图像学习细粒度评分建立了基础。本文沿用该任务表示与评分损失，以 SFibAI 为直接比较基线，进一步研究解剖切面和局部区域信息在评分过程中的组织与使用。

从整体图像特征得到一个准确分值，并不直接说明模型识别了哪个扫查切面、哪些区域呈现相关征象，以及这些区域怎样参与最终判断。全局表征可以包含上述信息，但仅输出评分时，使用者难以分别检查切面判断与区域响应。扫查切面和局部病灶标注为此提供了两类互补的监督：前者描述图像的解剖背景，后者指出值得关注的位置。将这些信息用于分级，需要在解剖背景、局部特征与图像级判断之间建立明确的计算联系；同时，还要适应局部标注并未覆盖完整病灶的监督条件。

多任务学习可以通过共享表征联合学习评分与空间任务 \citep{caruana1997multitask}，但共享表征本身并未规定预测的切面和定位结果如何进入评分计算。Grad-CAM 等事后归因方法能够突出与预测相关的图像区域 \citep{selvaraju2017gradcam}，其生成可视化的过程也不改变模型原有的评分路径。由此，本文关注的核心问题是：\emph{能否将切面与局部位置组织为直接参与预测的区域证据，在改善细粒度分级的同时，保留可供人工检查的空间输出？} 我们据此采用的设计原则是，在空间聚合之前引入切面条件，使局部特征先在相应的解剖背景下形成评分证据，再与整体图像判断结合。

为实现这一设计，本文提出 VCRE-Fib（View-Conditioned Regional Evidence for Fibrosis Grading），整体框架见图~\ref{fig:architecture}。该框架在保留全局评分路径的同时，利用切面标签与局部病灶框学习切面预测和弱定位。模型在空间聚合前生成切面条件化的区域评分证据，通过弱定位引导区域聚合，将所得修正加入全局预测，最后依据预测切面概率融合各条件输出。这一结构明确连接了空间信息的学习、使用与呈现：切面预测决定条件输出的融合，定位响应参与区域证据的聚合，二者也作为可查看的结果保留下来。推理阶段仅需输入图像，即可同时获得纤维化评分、扫查切面与定位图。

我们在统一的数据划分与评价协议下，将 VCRE-Fib 与 SFibAI 及两个独立训练的模块删除变体进行比较。在由 4 个中心、240 名患者和 4,107 张图像组成的测试集上，完整模型在本次比较中取得最低的预设综合风险（$\risk=0.186865$），相对 SFibAI 降低 7.115\%，并在图像级、patient-max、patient-median 及中心平衡 COR 上均取得最低值。模块删除比较同时显示，分级风险的降低并不伴随各项空间任务指标的一致提升，说明联合系统需要同时评价评分表现与空间输出质量。上述结果支持在当前任务中进一步研究切面条件区域证据的组织方式，并为连接细粒度评分与可检查的空间信息提供了具体框架；其对医生复核和临床决策的实际帮助仍需专门评价。
